\section{创业生存线：补贴战、下架、4000 块现金}

上一章讲产品，接下来进入创业过程，因为 AI 应用的产品判断和生存压力是绑在一起的。视频描述里提到补贴战争、产品下架、账上只剩 4000 块现金、怎么都融不到资。访谈中的情绪之所以强，是因为 Lovart 的爆发不是平滑增长，而是从连续挫折后突然出现的转折。

\begin{figure}[H]
\centering
\includegraphics[width=0.96\textwidth]{startup-survival-path.png}
\caption{应用创业生存线：补贴战、下架、现金见底、融资困难之后仍然活下来。自制概念图，依据视频描述和 00:36:00--00:45:00 对谈内容整理。}
\end{figure}

\begin{knowledgebox}{读图：低谷不是花絮，而是产品判断的一部分}
补贴战让团队看到流量不是壁垒，下架让团队意识到平台风险，现金见底迫使团队聚焦，融不到资考验创始人信念。所有这些都会影响产品路线。
\end{knowledgebox}

\begin{knowledgebox}{AI 应用创业风险表}
\begin{tabularx}{\textwidth}{p{0.22\textwidth}p{0.34\textwidth}X}
\toprule
风险 & 表现 & 对策 \\
\midrule
补贴战 & 靠钱买用户，留存不一定真实 & 尽快找到自然传播和付费场景。 \\
平台下架 & 外部平台规则改变或被迫下线 & 保留自有入口和用户关系。 \\
现金见底 & 账上资金不足以支撑试错 & 聚焦高价值功能，延长 runway。 \\
融资困难 & 市场不相信应用窗口 & 用真实用户和收入证明 PMF。 \\
\bottomrule
\end{tabularx}
\end{knowledgebox}

\begin{knowledgebox}{低谷对产品路线的影响}
\begin{tabularx}{\textwidth}{p{0.22\textwidth}p{0.34\textwidth}X}
\toprule
低谷事件 & 它迫使团队看见什么 & 可能带来的正确动作 \\
\midrule
补贴战 & 买来的增长不可靠 & 回到自然需求和真实留存。 \\
下架 & 平台入口不可控 & 建自有渠道和多平台备份。 \\
现金见底 & 试错时间有限 & 聚焦最强场景和最短变现路径。 \\
融不到资 & 叙事无法说服资本 & 用用户和收入反过来证明。 \\
\bottomrule
\end{tabularx}
\end{knowledgebox}

\subsection{战斗型 CEO}

前面讲外部压力，本节看这种压力如何落到创始人身上。视频描述说，陈冕更像一名战斗型 CEO。这个说法不是性格标签，而是 AI 应用创业的现实：产品、融资、增长、交付、团队情绪都压在同一个人身上。模型能力变化快，窗口期短，创始人必须在压力下快速判断。

\begin{figure}[H]
\centering
\includegraphics[width=0.90\textwidth]{fighting-ceo.png}
\caption{战斗型 CEO：产品、融资、增长、交付和情绪能量同时拉满。自制概念图，依据 00:02:00--00:03:10 与 01:19:00--01:28:00 对谈内容整理。}
\end{figure}

\begin{warningbox}{情绪高能不是管理系统}
创始人的兴奋和战斗力可以穿过低谷，但公司最终仍需要流程、组织和财务纪律。快乐可以提供能量，但不能替代管理。
\end{warningbox}

\subsection{本章小结}

Lovart 的创业故事说明，AI 应用窗口很短，低谷很硬。活到窗口打开，本身就是能力的一部分。

